\section{Exact bilinear decompositions}
\label{sec:bilinear-bounds}

For \(f:G\to\mathbb C\), write
\[
 \overline f=\frac1{q-1}\sum_{x\in G}f(x),
 \qquad
 f^\circ=f-\overline f\,\one_G.
\]

\begin{theorem}[Principal/nonprincipal identity]
\label{thm:principal-nonprincipal}
For arbitrary \(a,b:G\to\mathbb C\),
\begin{align}
 \mathcal I_H(a,b)
 ={}&
 \frac{2H}{q-1}
 \left(\sum_\omega a(\omega)\right)
 \left(\sum_r b(r)\right)
 +\mathcal R_H(a,b),
\label{eq:principal-decomposition}\\
 \mathcal R_H(a,b)
 ={}&
 \frac1{q-1}
 \sum_{\chi\ne\chi_0}
 \widehat{\one_{E_H}}(\chi)
 \widehat a(\overline\chi)
 \widehat b(\overline\chi).
\label{eq:nonprincipal-exact}
\end{align}
Only even nonprincipal characters contribute to
\eqref{eq:nonprincipal-exact}.
\end{theorem}

\begin{proof}
Insert multiplicative Fourier inversion for
\(\one_{E_H}(r\omega)\) into
\eqref{eq:incidence-form}.  The principal term is the first term of
\eqref{eq:principal-decomposition}.  The parity assertion follows
from \cref{thm:parity-nullspace}.
\end{proof}

\begin{corollary}[Spectral norm bound]
\label{cor:spectral-bound}
Let
\[
 M_H=\max_{\chi\ne\chi_0}
 |\widehat{\one_{E_H}}(\chi)|.
\]
Then
\begin{equation}
 |\mathcal R_H(a,b)|
 \le
 M_H\|a^\circ\|_2\|b^\circ\|_2
 \ll
 \sqrt q\log q\,
 \|a^\circ\|_2\|b^\circ\|_2.
\label{eq:spectral-bilinear-bound}
\end{equation}
\end{corollary}

\begin{proof}
Use Cauchy--Schwarz in \eqref{eq:nonprincipal-exact}, followed by
Parseval and \cref{cor:PV-bound}.  Odd characters already vanish.
\end{proof}

\begin{proposition}[Exact centered operator norm]
\label{prop:centered-operator}
\[
 \left\|
 K_Hb-\frac{2H}{q-1}
 \left(\sum_rb(r)\right)\one_G
 \right\|_2
 \le M_H\|b^\circ\|_2.
\]
The best constant is exactly \(M_H\).
\end{proposition}

\begin{proof}
On the character basis, the centered operator deletes the principal
mode and has the nonprincipal singular values from
\cref{thm:complete-spectrum}.  Their maximum is \(M_H\), and an
extremizing character gives equality.
\end{proof}

\begin{remark}[Two distinct tasks]
The principal term in \eqref{eq:principal-decomposition} is a density
term of size \(2H/(q-1)\).  The remainder is a multiplicative
dispersion term.  A proof must control both; a small nonprincipal
spectrum cannot compensate for an unaffordable principal mass.
\end{remark}
